\subsection{Завершение глобальной классификации равенства}\label{subsec:complete10equality}

Строгий расчёт доказывает минимум, но сам по себе не решает задачу равенства. Переход от $D\le54672$ к $D=54675$ меняет несколько промежуточных окон определителей и может восстановить ветви, ранее удалённые неравенством подъёма. Следующий дополнительный аргумент проверяет оба изменения.

\begin{theorem}\label{thm:global10equality}
Пусть $L$ --- положительно определённая чётная целая решётка ранга 10 с $m(L)\ge6$ и $d(L)=54675$. Тогда $L$ содержит подрешётку, изометричную $\sqrt3E_8$, и
\[
L\cong K_{10}\quad\text{или}\quad L\cong\Kalt.
\]
Эти две решётки неизометричны. Таким образом, $e_{10}(6)=54675$ имеет ровно два глобальных класса равенства. Внешние и вычислительные предпосылки совпадают с указанными в теореме \ref{thm:e10exact}, с добавлением полных сертификатов границы равенства, описанных ниже.
\end{theorem}

\begin{proof}
Предположим, что $L$ не содержит $\sqrt3E_8$. Выберем последовательные примитивные кратчайшие двойственные сечения, как в \eqref{eq:dualflag}. Верхние окна при равенстве, округлённые с помощью точных целых степеней и сравнений определителей, приведены ниже. Все нижние границы остаются прежними.
\begin{center}
\begin{tabular}{crrr}
\toprule
Ранг & Строгая граница & Граница равенства & Новый определитель\\\midrule
1&6&6&---\\
2&44&44&---\\
3&210&210&---\\
4&789&789&---\\
5&2492&2492&---\\
6&6459&6460&6460\\
7&13938&13940&13940\\
8&24673&24676&24676\\
9&40054&40056&40056\\
10&54672&54675&54675\\
\bottomrule
\end{tabular}
\end{center}
В частности, нижняя граница ранга 8 остаётся равной 6564: определитель 6561 вынуждал бы наличие $\sqrt3E_8$ вопреки предположению. Ранги до пятого включительно не меняются.

\textbf{Восстановление ранее исключённых ветвей.}
Отсечение продолжения имеет вид
\[
3d_r^2>4d_{r-1}U_{r+1}.
\]
Увеличение $U_{r+1}$ может нарушить это неравенство, даже если $d_r$ находится в старом интервале. Поэтому недостаточно проверить новые граничные определители над прежними сохранёнными списками. Каждое прежнее отсечение проверено заново при увеличенных верхних границах:
\begin{center}
\begin{tabular}{P{0.50\linewidth}rr}
\toprule
Семейство & Старые отсечения & Восстановлено\\\midrule
Полные поиски от родителей ранга 6 &3252&0\\
Исходные записи продолжений ранга 7 &52208&4\\
Исходные записи продолжений ранга 8 &18110&0\\
Исходные записи продолжений ранга 9 &3819&0\\
\bottomrule
\end{tabular}
\end{center}
Четыре восстановленные матрицы ранга 7 имеют $d_6=4131$ и $d_7=11658$. Их следующий определитель удовлетворяет
\begin{equation}\label{eq:restored7floor}
d_8\ge\frac{3\cdot11658^2}{4\cdot4131}
=\frac{3775249}{153},\qquad
24673<\frac{3775249}{153}<24676.
\end{equation}
В расширенном окне сравнения определителя ранга 8 поэтому вынуждают $d_8=24676$. Все четыре исходные матрицы сохранены для дополнительной проверки; новое объединение по изометриям не требуется.

\textbf{Независимый сертификат необходимых условий.}
Для границы равенства используется второй, более простой расчёт. Он не импортирует исходный генератор расширений или его перечислитель векторов и не требует минимумов дискриминантных классов. Пусть $G$ --- родительская матрица ранга $r$, $d=\det G$, $A=\operatorname{adj}G$. Целое расширение одним вектором с чётной последней диагональю $n$ и столбцом скалярных произведений $p$ имеет целевой определитель
\[
T=nd-p^{\mathsf T}Ap.
\]
Его класс кодируется величинами
\[
c=Ap\pmod d,\qquad Q(c)=p^{\mathsf T}Ap\pmod{2d}.
\]
Значение $Q$ определено корректно: замена $p$ на $p+Gz$ меняет квадратичное выражение на $2d\,z^{\mathsf T}p+d\,z^{\mathsf T}Gz$, что кратно $2d$ благодаря чётности $G$. Ядро отображения $p\mapsto Ap\pmod d$ равно $G\Z^r$, поэтому классов ровно $d$.

Классы и значения строятся от $c=0$, $Q(0)=0$ добавлением столбцов $Ae_j$ по точной рекурсии
\begin{equation}\label{eq:finitequadraticrecurrence}
Q(c+Ae_j)=Q(c)+2c_j+A_{jj}\pmod{2d}.
\end{equation}
Программа проверяет согласованность при каждом повторном достижении класса и устанавливает, что достигнуто ровно $d$ классов. Чётность $n$ требует
\begin{equation}\label{eq:boundaryquadraticcondition}
Q(c)\equiv-T\pmod{2d}.
\end{equation}
Допустима также предварительная проверка общего делителя: если
\[
g=\gcd\bigl(2d,\ (A_{ii})_i,\ (2A_{ij})_{i<j}\bigr)
\]
не делит $T$, подходящих классов нет.

Для каждого класса, удовлетворяющего \eqref{eq:boundaryquadraticcondition}, новая перпендикулярная нормаль должна быть кратчайшей в двойственной решётке. Согласно \eqref{eq:integerdualtest}, необходимо
\begin{equation}\label{eq:boundarydualcondition}
T(q^{\mathsf T}Aq)+\operatorname{dist}(q^{\mathsf T}c,d\Z)^2\ge d^2
\end{equation}
для каждого ненулевого целого $q$ с $T(q^{\mathsf T}Aq)<d^2$. Независимый точный рациональный перечислитель строит все требуемые $q$ внутри $q^{\mathsf T}G^{-1}q\le d/T$. Один нарушающий условие вектор $q$ исключает класс. Если классы всё ещё остались, расширение, которое может продолжаться следующим членом полного кратчайшего двойственного флага, обязано удовлетворять следующему неравенству подъёма
\begin{equation}\label{eq:boundarynextlift}
3T^2\le4d\,U_{r+2}.
\end{equation}
Все классы, пережившие предыдущие проверки в этом расчёте, нарушают \eqref{eq:boundarynextlift}. Не проверять полный минимум расширения безопасно: используются необходимые условия, уже исключающие все ветви.

Полное дополнительное покрытие состоит из 4952 пар «родитель--целевой определитель»:
\begin{center}
\begin{tabular}{P{0.33\linewidth}P{0.10\linewidth}P{0.11\linewidth}P{0.18\linewidth}P{0.14\linewidth}}
\toprule
Источник родителей & Число & Цель & Осталось классов & Проходит\\\midrule
Прежние представители ранга 5 &925&6460&58&0\\
Прежние представители ранга 6 &3752&13940&20&0\\
Прежние представители ранга 7 &224&24676&0&0\\
Восстановленные записи ранга 7 &4&24676&0&0\\
Прежние представители ранга 8 &47&40056&0&0\\
\bottomrule
\end{tabular}
\end{center}
Столбец «осталось классов» считает классы после сравнения определителя и проверки кратчайшей нормали, до следующего неравенства подъёма. Это вычислительные записи, а не количества типов изометрии. 58 классов возникают у четырёх родительских матриц, а 20 --- у шести. Все нарушают следующее неравенство подъёма. Все четыре восстановленные формы ранга 7 исключаются непредставимостью целевого остатка в их конечных квадратичных модулях.

Решающие сравнения следующего шага можно видеть непосредственно. Определители родителей для 58 классов ранга 6 равны $1934,1830,1672,1600$, с кратностями $2,4,4,48$. Для 20 классов ранга 7 они равны $4428,4443,4268,4379,4355,4920$, с кратностями $2,4,4,4,4,2$. Даже наибольший родительский определитель в каждом семействе даёт
\[
\begin{aligned}
3\cdot6460^2&=125194800>4\cdot1934\cdot13940=107839840,\\
3\cdot13940^2&=582970800>4\cdot4920\cdot24676=485623680.
\end{aligned}
\]
Таким образом, все 78 оставшихся классов явно исключены из полного флага.

Для дополнительного контроля полноты 4952 пары разбиты на непересекающиеся категории исключения: 3960 удалены известным двойственным вектором родителя, 493 --- проверкой общего делителя, 420 --- непредставимостью в конечном модуле, 69 --- явными более короткими двойственными векторами для каждого подходящего класса, 10 --- следующим неравенством подъёма для оставшихся классов. Сумма равна ровно 4952.

\textbf{Почему покрыто всё окно равенства.}
Префикс до ранга 5 включительно не изменился. Любая новая форма ранга 6 возникала бы при определителе 6460 у одного из 925 прежних родителей, но ни одна не продолжается в ранг 7. Следовательно, в ранге 7 новые возможности --- только конец 13940 над прежними родителями ранга 6 и четыре восстановленные записи. Случаи нового конца не продолжаются. В ранге 8 конец 24676 проверен над всеми 224 прежними представителями и всеми четырьмя восстановленными записями; расширений не осталось. Значит, нет нового родителя ранга 8, который пришлось бы продолжать по остальной части интервала ранга 9. Конец 40056 над всеми 47 прежними представителями ранга 8 также исключён. Старый низкий интервал ранга 9 остаётся пустым, а каждая запись старого верхнего интервала по-прежнему удовлетворяет
\[
d_{10}\ge\frac{12966723}{212}>54675.
\]
Таким образом, исключены все прежние и вновь допустимые флаги. Аргумент использует ранее сертифицированные переходы к представителям и полную проверку восстановленных отсечений; он не предполагает, что гипотетический минимизатор содержит $K_9$ или имеет базис из минимальных векторов.

Получено противоречие существованию случая равенства без $\sqrt3E_8$. Поэтому каждый минимизатор содержит эту подрешётку. Теорема \ref{thm:rank10e8} оставляет ровно две возможности $K_{10}$ и $\Kalt$; ранги их матриц по модулю 3 равны соответственно 4 и 3, так что они неизометричны и как абстрактные целые решётки. Теорема доказана.
\end{proof}

Независимые проверки равенства находятся в файлах \file{research_20261008/independent_review/check_equality_windows_and_recovered_prunes.py} и \file{research_20261008/independent_review/check_boundaries_by_quadratic_module.py}. Статусы сертификатов --- \file{PASS} и \file{ALL_EXCLUDED}; вторая проверка охватывает все 4952 пары. Эти файлы используют ранее независимо реализованную рациональную матричную и векторную арифметику, а не исходное ядро расширений. Полный сертификат прежнего строгого поиска остаётся зависимостью. Отдельное публичное доказательство $\gamma_9\le2$ по-прежнему служит внешней предпосылкой. Ни одна из этих внутренних проверок не является внешним рецензированием.

Для дополнительного сопоставления с исходным методом расширений заново выполнены все 5010 родительских задач: 4952 граничные пары и 58 попыток продолжения новых матриц ранга 6. Из них 4018 сразу исключены родительскими оценками подъёма, а в 992 полностью перечислены таблицы дискриминантных классов. Получено ровно 58 исходных матриц ранга 6 и 20 ранга 7. Ни одна не входит в полный флаг, а добавления рангов 8 и 9 пусты. Отдельный независимый аудит проверяет настройки каждой задачи и связь с родителем, независимо проверяет 78 исходных матриц, их прямые и двойственные минимумы и сопоставляет точные отпечатки классов с сертификатом конечных модулей. Новая генерация и независимый аудит выполнены; второй полный повтор генерации доступен как опция, но не заявляется как ещё раз выполненный.

Основная новая проверка --- \file{research_20261008/rank10_equality/verify_equality_certificate.py}; сохранён статус \file{PASS_STATIC_EQUALITY_CERTIFICATE} с маркером \file{GLOBAL_RANK10_EQUALITY_TWO_TYPES_PASS}. Независимый аудит задач --- \file{research_20261008/independent_review/check_delta_chain_independent.py}. После проверки базового сертификата основной файл допускает \texttt{--skip-base}; необязательный параметр \texttt{--replay} заново генерирует добавленные задачи. Статическая проверка сертификата и полный повтор генерации являются разными операциями, их статусы фиксируются отдельно.
